\documentclass[11pt,letterpaper]{article}
\usepackage[T1]{fontenc}
\usepackage{lmodern,amsmath,amssymb,amsthm,microtype,xcolor}
\usepackage[left=1.05in,right=1.05in,top=0.95in,bottom=0.95in]{geometry}
\usepackage{hyperref}
\definecolor{inkblue}{RGB}{25,65,100}
\hypersetup{colorlinks=true,allcolors=inkblue,pdftitle={VP versus VNP by separating tensor orbit closures}}
\newtheorem{theorem}{Theorem}
\newtheorem{lemma}[theorem]{Lemma}
\newtheorem{corollary}[theorem]{Corollary}
\newcommand{\GL}{\mathrm{GL}}
\newcommand{\Cl}{\operatorname{Clique}}
\newcommand{\bal}{\mathcal B}
\newcommand{\rk}{\operatorname{rank}}
\title{VP versus VNP by separating tensor orbit closures}
\author{}\date{November 27, 2026}
\AddToHook{cmd/maketitle/after}{%
  \par\noindent\begin{minipage}{\linewidth}
  \small\sffamily
  I asked ChatGPT these questions to understand what follows from the recent
  VP versus VNP breakthrough. I share its answers so that these explorations
  are available beyond a private conversation. These are AI-generated notes,
  not research papers. I do not claim authorship or priority for the ideas.
  Arguments may be wrong, and earlier work may be overlooked or misattributed.
  I welcome corrections by mail or via
  \href{https://github.com/radomite/algebraic-complexity/issues}{GitHub}.
  \end{minipage}\par\medskip
}
\AddToHook{shipout/foreground}{%
  \ifnum\ReadonlyShipoutCounter>1\relax
  \put(0pt,-24pt){\makebox[\paperwidth][c]{%
    \normalfont\sffamily\fontsize{8}{10}\selectfont\color[gray]{0.38}%
    AI-generated exploration notes. Arguments and attributions may be wrong.
    \href{https://github.com/radomite/algebraic-complexity/issues}{Corrections welcome on GitHub.}}}%
  \fi
}

\begin{document}
\maketitle

\section{The polynomial and its orbit-closure witness}
Work with formal polynomials over a field $F$. Size counts the gates of
a binary arithmetic circuit with variables and field constants, allowing
reuse of every gate. For a host with $N$ vertices set
\[
 \Cl_{N,k}=\sum_{\substack{U\subseteq[N]\\|U|=k}}
      \prod_{i\in U}X_i\prod_{\{i,j\}\subseteq U}E_{ij},
 \qquad \Cl_N=\sum_{k=0}^N\Cl_{N,k}.
\]
We recover the clique bound of the wire argument in
\cite[Appendix E]{anthropic}. All orbit closures below are taken over
$K=\overline F$. Extending scalars preserves every circuit over $F$.

\begin{theorem}\label{main}
For each fixed $k\ge2$ the circuit size $C_F$ satisfies
\[
 \liminf_{N\to\infty}\frac{\log C_F(\Cl_{N,k})}{\log N}
                  \ \ge\ \frac{k^{3/2}}{3(k-1)} .
\]
The witness is a noncontainment of tensor orbits: for $q=2^\ell$ and
every fixed $\epsilon>0$, a clique tensor lies outside the orbit closure
for rank budget $r$ whenever $1\le r\le q^{k^{3/2}/3-\epsilon}$ and
$\ell$ is sufficiently large. The separating equation can be chosen
over the prime field, of degree at most $q^{k(k-1)}$.
\end{theorem}

Give each edge of $K_k$ an independent label in $[q]$. A label assignment
$\sigma$ supplies vertex $i$ with its incident tuple $\sigma_i\in[q]^{k-1}$.
Consider the polynomial $\sum_\sigma\prod_i x_{i,\sigma_i}$.
It is a projection of $\Cl_{N,k}$ at $N=kq^{k-1}$. Make one host palette
for each $i$, with one vertex of weight $x_{i,\rho}$ for each tuple $\rho$.
Give intrapalette edges weight zero. Across palettes $i,j$, set the edge
weight to one exactly when the two tuples agree on label $ij$, and to
zero otherwise. Every surviving clique picks one vertex per palette and
each consistent choice is counted once, in every characteristic.

Let $W_i=(K^q)^{\otimes(k-1)}$ have its tuple basis and
$W=\bigotimes_i W_i$. The coefficient of one variable from each palette
is an entry of a tensor in $W$. For our projected polynomial the tensor is
\[
               T=\mathop{\boxtimes}\limits_{\{i,j\}\in E(K_k)}
                        \sum_{a=1}^q e_a^{(i,j)}\otimes e_a^{(j,i)} .
\]
Here $(i,j)$ denotes the end of edge $ij$ owned by $i$, and $\boxtimes$
groups factors at the same site. The displayed factor is an equality
wire. On a triangle, for instance, the monomials have the form
$x_{1,(a,b)}x_{2,(a,c)}x_{3,(b,c)}$: the three wires carry $a,b,c$.

\section{One source orbit for all balanced cuts}
Write $W_S=\bigotimes_{i\in S}W_i$. A cut $S\mid S^c$ is balanced if
$k/3\le |S|\le2k/3$. Take one representative per unordered cut in
$\bal$. For $r\ge1$ let $\Phi_r$ take arbitrary pairs of tables to their sum
\[
    \Phi_r(a,b)=\sum_{S\in\bal}\sum_{j=1}^r a_{S,j}\otimes b_{S,j},
       \qquad a_{S,j}\in W_S,\quad b_{S,j}\in W_{S^c}.
       \tag{1}\label{join}
\]
For one cut the inner sum is a matrix of rank at most $r$, with rows
indexed by $S$ and columns by $S^c$. Write $J_r=\overline{\mathrm{im}\,\Phi_r}$.
It is the affine join of these matrix rank varieties.
For a tensor on several spaces, its orbit consists of its images under
independent invertible linear maps on those spaces. Its Zariski closure
is the common zero set of all polynomials in the entries that vanish on
the orbit. Thus one polynomial that vanishes there and is nonzero at the
clique suffices to distinguish the two.

\begin{lemma}\label{cut}
The coefficient tensor of a size-$s$ circuit on the $k$ palettes lies
in $\mathrm{im}\,\Phi_{2s}$.
\end{lemma}
\begin{proof}
At gate $g$, denote by $[g]_U$ the component using one variable from
each palette in $U$ and none from the others. At a product,
$[gh]_U=\sum_{A\sqcup B=U}[g]_A[h]_B$ with ordered splits. Expand from
the output on $[k]$. Expand both summands at an addition. At a product
split follow the larger child (ties go left), saving the other factor.
Stop the first time the followed set has size at most $2k/3$. Its parent
had size greater than $2k/3$, so the stopping set $S$ has size greater
than $k/3$. An input with a larger-than-allowed support contributes zero.
For each stopping pair $(g,S)$ the factor is always $[g]_S$.
All saved contexts use exactly the complementary palettes and can be
added together. Even if a gate occurs along many paths this gives one
product per pair. There are at most $s$ gates for each oriented side and
two orientations per cut.
\end{proof}

\begin{lemma}\label{orbit}
There is an explicit tensor $Q_r$ on a larger space, with product maps
$\iota,\pi$ satisfying $\pi\iota=1_W$, for which
\[
                 A\in J_r\quad\Longleftrightarrow\quad
                 \iota A\in\overline{G\cdot Q_r}.
\]
Here $G$ is the product of the general linear groups of the source factors.
\end{lemma}
\begin{proof}
Keep one terminal per ordered edge end and add two private hubs for each
cut. For a fixed cut $S$, connect every terminal by a $q$-wire to the hub
on its side, and connect the two hubs by an $r$-wire. Let $Q_S$ be the
product of these wires, grouping all indices incident at each party.
Other hubs have dimension one in this tensor. Form $Q_r=\bigoplus_S Q_S$
by taking the direct sum of the local spaces at every party and putting
each $Q_S$ in its own diagonal block.

Apply arbitrary linear maps from terminal spaces to $K^q$ and from
hub spaces to $K$. These restrictions of $Q_r$ give exactly
\eqref{join}. Indeed, the middle index selects $j$, and each hub
functional supplies a table on all end labels on its side and $j$.
The terminal maps can be absorbed in those tables. Conversely choosing
identity terminal maps on each block and prescribing the hub entries
realizes every list of tables independently.

Embed these output spaces factorwise into coordinate subspaces of the
corresponding source spaces by $\iota$, and use the coordinate
projections as a factorwise left inverse $\pi$.
Composing any restriction with $\iota$ gives a product endomorphism.
Such maps lie in the closure of $G$ since invertible matrices are dense.
Thus $\iota(\mathrm{im}\,\Phi_r)\subseteq\overline{G\cdot Q_r}$.
Conversely, $\pi$ takes every orbit point to a restriction, and hence
takes its closure into $J_r$. Taking closures gives both implications.
\end{proof}

\section{An equation separating the two tensors}
A spectral point $\phi$ assigns a positive real number to every nonzero
tensor on a fixed finite site set, and zero to the zero tensor. It is
multiplicative on grouped tensor products, additive on block direct sums,
and monotone under restrictions. In particular
$\phi(\sum_{j=1}^L A_j)\le L\max_j\phi(A_j)$: the sum is a restriction
of the direct sum. A nonzero scalar factor preserves $\phi$ by invertible
rescaling. We use the following result from \cite{euclid}.

\begin{quote}
\textbf{Euclidean binary-wire theorem.}
For any fixed finite Euclidean site configuration with at least two
distinct sites there are a spectral point $\phi$ over $K$ and $c>0$
such that $c\log_2\phi(\text{binary wire from }u\text{ to }v)=\|u-v\|$
simultaneously for all pairs.
\end{quote}

Write $H=c\log_2\phi$, with $H(0)=-\infty$. A $2^\ell$-wire costs
$\ell$ times its length. The following elementary observation lets this
exact-restriction theorem produce an equation on a Zariski closure.

\begin{lemma}\label{equation}
Let $X$ be a set of tensors in a $v$-dimensional tensor space, all on
the same sites as $T$. If $H(T)>\sup_X H+2c$, there is a homogeneous
polynomial $E$ of degree $v$ with $E(X)=0$ and $E(T)=1$.
\end{lemma}
\begin{proof}
Set $d=v$. Up to regrouping, the coordinates of $A^{\boxtimes d}$ are
monomials of degree $d$ in $v$ entries. Thus the span of all such powers
has dimension at most $\binom{v+d-1}{d}<4^d$. If $T^{\boxtimes d}$ lay
in the span of powers from $X$, it would be a sum of fewer than $4^d$
scaled powers from $X$, by choosing a basis from the spanning set.
The sum inequality and multiplicativity would give
$dH(T)\le d\sup_X H+2cd$. A linear functional zero on that span and
one at $T^{\boxtimes d}$ therefore exists. Evaluate it on
$A^{\boxtimes d}$ to obtain $E(A)$. The argument uses ordinary tensor
coordinates and works in positive characteristic as well.
\end{proof}

We now price $X=\mathrm{im}\,\Phi_r$. Fix $k$ and put $M=\binom k2$.
Place terminal $(i,j)$ at $u_{ij}=(e_i-e_j)/\sqrt2$ in $\mathbb R^k$.
The two ends of each target wire are distance two apart, giving
$H(T)=2M\ell$. The vector sum on the $S$ side is
\[
 v_S=\frac{k{\bf1}_S-|S|{\bf1}}{\sqrt2},
 \qquad \|v_S\|^2=\frac{k|S|(k-|S|)}2,\qquad
 b:=\min_{S\in\bal}\|v_S\|\ge k^{3/2}/3 .
\]
Indeed the two ends of each internal edge cancel; the last inequality
uses balance. Fix $0<t<1$. Place the two hubs for $S$ at
$\pm t v_S/\|v_S\|$, the positive one serving $S$. Apply the import
once to the terminals and all these hubs. The scale $c=c_{k,t}$ is
independent of $q,r$. At unused hubs adjoin scalar factors to the output.

For unit vectors $u,w$, the tangent bound for the square root gives
$\|u-tw\|\le1-t\langle u,w\rangle+t^2/2$.
Since the terminal sums on the two sides are $v_S$ and $-v_S$, the
total length of the $2M$ spokes is at most $2M-2tb+Mt^2$.
The middle wire costs at most $2t\lceil\log_2 r\rceil$, by restricting
a power-of-two wire. Sum the restrictions of the cut networks on the
common sites. The resulting inequality is
\[
 H(T)-\sup_{A\in X}H(A)
 \ \ge\ 2t\bigl(b\ell-\lceil\log_2 r\rceil\bigr)
                -Mt^2\ell-c\log_2|\bal| .
 \tag{2}\label{gap}
\]
If $r\le q^{b-\epsilon}$, choose $0<t<\min(1,\epsilon/M)$.
The right side grows linearly in $\ell$ and eventually exceeds $2c$.
Lemma~\ref{equation} supplies $E$ on $W$, where $v=q^{k(k-1)}$.
It vanishes on $J_r$ because its zero set is closed. In the common
ambient space of Lemma~\ref{orbit}, $E\circ\pi$ vanishes on
$\overline{G\cdot Q_r}$ and takes value one at $\iota T$.
The translates of this equation span a $G$-stable space of equations.
This is the promised orbit-closure separation.

The equation also has a finite coefficient prescription. Give a
degree-$v$ polynomial $E$ unknown coefficients, equate all coefficients
of $E(\Phi_r(a,b))$ to zero in the independent table parameters, and
impose $E(T)=1$. These linear equations have entries in the prime
field. Consistency over $K$ implies consistency there by row reduction.
This specifies a separator at each stated budget without an assumption
on the cost of evaluating it.

\section{From noncontainment to the two class separations}
At $N_\ell=k2^{(k-1)\ell}$, Lemma~\ref{cut} and the clique projection
put $T$ in $J_{2s}$ whenever $\Cl_{N_\ell,k}$ has size $s$.
Inequality~\eqref{gap} and the separating equation exclude this for
$2s\le 2^{(b-\epsilon)\ell}$ and all sufficiently large $\ell$.
For $N_\ell\le N<N_{\ell+1}$ zero the extra vertex weights. This gives
the same lower bound and $\log_2 N<\log_2 k+(k-1)(\ell+1)$.
Dividing and letting $\epsilon$ decrease to zero proves Theorem~\ref{main}.

\begin{corollary}
Over every field $F$, $\mathrm{VFPT}_F\subsetneq\mathrm{VW}[1]_F$
and $\mathrm{VP}_F\subsetneq\mathrm{VNP}_F$.
\end{corollary}
\begin{proof}
An FPT circuit family has size $f(k)N^a$ for a common exponent $a$.
Choose $k$ with $k^{3/2}/(3(k-1))>a$ and then let $N$ grow.
The fixed factor $f(k)$ disappears in the exponent, contradicting
Theorem~\ref{main}.

For membership, take Boolean selectors $y_i$ and form
\[
  \prod_i(1+y_i(X_i-1))
      \prod_{i<j}(1+y_i y_j(E_{ij}-1)).
\]
Each assignment contributes precisely the weight of its selected set.
Summing over assignments of Hamming weight $k$ gives $\Cl_{N,k}$.
Each factor has constant bounded-fan-in depth, with a single unbounded
product above them. This is a weight-$k$ sum of a weft-one circuit as
in the definition of $\mathrm{VW}[1]$ \cite{be}. That class contains
$\mathrm{VFPT}$ by its allowed fixed-parameter substitutions: the sum
of $z$ over the weight-one assignments to a single selector equals $z$.

Summing over all selectors gives $\Cl_N$ as a Boolean sum of a
polynomial-size, polynomial-degree circuit, hence a VNP family.
If it had polynomial-size circuits, replace each $X_i$ by $uX_i$ and
keep the coefficients of $u^0,\ldots,u^k$ at each gate. Truncated
convolution costs $O((k+1)^2)$ per product. The $u^k$ coefficient is
$\Cl_{N,k}$, which would give the forbidden FPT circuit family.
\end{proof}

\bibliographystyle{../linkedtitles}
\bibliography{../references}
\end{document}
